% Generated draft: requires mathematical and source review.
\documentclass[11pt]{article}
\usepackage[margin=25mm]{geometry}
\usepackage{fontspec}
\setmainfont{Latin Modern Roman}
\usepackage{amsmath,amssymb,mathtools,unicode-math}
\setmathfont{Latin Modern Math}
\usepackage{enumitem,needspace,xurl,hyperref}
\setlength{\parindent}{0pt}
\setlength{\parskip}{5pt}
\setlength{\emergencystretch}{3em}
\newcommand{\N}{\mathbb{N}}
\newcommand{\Z}{\mathbb{Z}}
\newcommand{\Q}{\mathbb{Q}}
\newcommand{\R}{\mathbb{R}}
\newcommand{\C}{\mathbb{C}}
\newcommand{\F}{\mathbb{F}}
\newcommand{\A}{\mathbb{A}}
\newcommand{\PP}{\mathbb{P}}
\newcommand{\E}{\mathbb{E}}
\newcommand{\eps}{\varepsilon}
\newcommand{\dd}{\,\mathrm{d}}
\begin{document}
\section*{UnsolvedMath ID: 10400140}
\textbf{Source problem code:} catalogue entry
\subsection{Definitions and mathematical statement}
Let \((W,L,\rho)\) be any triple admissible for the source's quantum hyperbolic invariants (QHI). Let \(K_N(W,L,\rho)\) denote its invariant at the admissible positive integer parameter \(N\), and let \(c_I(W,L,\rho)\) and \(R\) be the source's class and regulator-type expression. The real volume conjecture attributed to S. Baseilhac and R. Benedetti is
\[\lim_{N\to\infty}\frac{2\pi}{N^2}\log\bigl|K_N(W,L,\rho)\bigr|
=\operatorname{Im}R\bigl(c_I(W,L,\rho)\bigr)\]
for every such triple. The logarithm is the real natural logarithm of the absolute value, and the denominator is \(N^2\), exactly as supplied. This is an ordinary limit along the parameters for which the source defines \(K_N\), not an optimistic limit or a logarithm with complex branch choices. Root-of-unity multiplication does not change the absolute value.

If \(V=\operatorname{Im}R(c_I(W,L,\rho))\) is a well-defined finite real number, the finite-limit reading means that for every \(\varepsilon>0\) there is \(N_0\) such that for every admissible \(N\geq N_0\), the logarithm is defined and
\[\left|\frac{2\pi}{N^2}\log|K_N(W,L,\rho)|-V\right|<\varepsilon.\]
The source does not specify how zero invariant values are treated; an extended convention \(\log0=-\infty\) would still require addressing them for a finite limit. The domain and equivalence data of the triples, allowed \(N\), definition of \(c_I\), domain and ambiguity of \(R\), and well-definedness of its imaginary part must be supplied by the surrounding theory. The formula does not assert that every \(W\) is hyperbolic or replace its right side by a particular geometric volume without those definitions.
\subsection{Short English statement}
The absolute values of the QHI invariants should have logarithmic growth \(N^2/(2\pi)\) times the imaginary regulator value, for every admissible triple. The omitted QHI definitions and possible zero values require clarification.
\subsection{Sources}
\begin{itemize}
\item Local UnsolvedMath registry, permanent ID \(10400140\), exact source code \texttt{catalogue entry}.
\item S. Baseilhac and R. Benedetti, Real Volume Conjecture for QHI, in Ohtsuki, \emph{Problems on invariants of knots and 3-manifolds} (2002), Conjecture 7.25, PDF page 116: \url{https://msp.org/gtm/2002/04/gtm-2002-04-024s.pdf}. Recorded access 2026-07-29; extraction pdf-font-structure.
\end{itemize}
\end{document}
